% source: manti.01.ma, 1ère SM, devoir surveillé n°1 (Drive 0B6pfSJcEVng0aGZrVTRjakZMNjg)
\ex{Fonctions}{aGZr}{Ex. 1}{7}
On considère les fonctions $f$ et $g$ définies par $f(x)=-x^2-2x+3$ et $g(x)=\dfrac{2(x-1)}{x+1}$.
\begin{enumerate}
\item Dresser le tableau de variation de chacune des fonctions $f$ et $g$.
\item Résoudre l'équation $f(x)=0$, puis donner une interprétation géométrique du résultat.
\item Tracer les courbes $(C_f)$ et $(C_g)$.
\item Déterminer graphiquement les solutions de l'inéquation $x^2+2x\leq\dfrac{x+5}{x+1}$.
\end{enumerate}
\fin
\ex{Fonctions}{aGZr}{Ex. 2}{8}
Soit $f$ la fonction définie par $f(x)=\dfrac{x^2-x}{x^2-x+2}$.
\begin{enumerate}
\item Déterminer l'ensemble de définition de $f$.
\item \begin{enumerate}
\item Montrer que la fonction $f$ est majorée par $1$.
\item La fonction $f$ admet-elle une valeur maximale ? Justifier.
\end{enumerate}
\item Montrer que $f$ admet une valeur minimale en $a=\dfrac12$.
\item On pose $g(x)=x^2-x$.
\begin{enumerate}
\item Déterminer une fonction de référence $h$ telle que $f=h\circ g$.
\item Étudier la monotonie de $f$ sur l'intervalle $\left]-\infty,\dfrac12\right]$.
\end{enumerate}
\end{enumerate}
\fin
\ex{Fonctions}{aGZr}{Ex. 3}{8}
On considère la fonction $f$ définie par $f(x)=\dfrac{\sqrt x}{x+1}$.
\begin{enumerate}
\item \begin{enumerate}
\item Déterminer $D_f$ et montrer que $\dfrac{f(x)-f(y)}{x-y}=\dfrac{1-\sqrt{xy}}{(x+1)(y+1)\left(\sqrt x+\sqrt y\right)}$.
\item Étudier la monotonie de $f$ sur chacun des intervalles $[0,1]$ et $[1,+\infty[$.
\end{enumerate}
\item Soient $p$ et $q$ deux entiers de $\mathbb{N}^*$. Montrer que $0<\dfrac{\sqrt{pq}}{(p+1)(q+1)}\leq\dfrac14$.
\end{enumerate}
\fin
